\section{Related Work}
\label{sec:related}

\paragraph{Neural knowledge graph embeddings.}
TransE~\cite{transe2013}, DistMult~\cite{distmult2015}, ComplEx~\cite{complex2016},
ConvE~\cite{conve2018}, and RotatE~\cite{rotate2019} established the standard
link-prediction setup on FB15K-237 and WN18RR.
None of these models expose a structural notion of consistency, and none
emit a derivation for predictions.

\paragraph{LLM-based KG reasoning.}
KG-BERT~\cite{yao2019kgbert} reframes link prediction as triple
classification with a pretrained language model. KGT5~\cite{saxena2022sequence}
casts it as sequence-to-sequence over verbalized triples, achieving
MRR~$0.300$ on FB15K-237. SimKGC~\cite{wang2022simkgc} adds contrastive
training over LM-encoded entity descriptions, reaching MRR~$0.336$. These
methods improve raw retrieval accuracy but provide no structural signal
when their predictions are inconsistent with the underlying ontology---a
gap our sheaf-cohomology layer directly addresses.

\paragraph{Sheaves and topology in machine learning.}
Sheaf neural networks~\cite{bodnar2022sheaf, hansen2020sheaf} and topological
data analysis have established the toolkit; we apply it to KG
inconsistency rather than message passing or persistence.

\paragraph{Categorical knowledge representation.}
Ologs~\cite{spivak2012olog} provide the syntactic substrate for our sheaf
construction. The broader Olog-as-category-of-types perspective belongs to
applied category theory~\cite{fongspivak2019}, where structure-preserving
functors between categories are the standard tool for relating
representations --- the same lens we apply in treating a KG as a sheaf
over an Olog.

\paragraph{Constrained decoding.}
Constrained decoding restricts the output distribution at sampling time
to a domain-specified admissible set. Relevant lines include
grammar-constrained generation~\cite{geng2023flexiblegrammar} and
schema-constrained generation~\cite{willard2023outlines}. Our (D) component
specializes this to Olog-derived admissible sets per generation step.
The novel contribution is the \emph{architectural separation} between (B)
information-flow constraints at attention and (D) output-validity
constraints at decoding, with the proof object as their shared interface.
